\section{Introduction}

Cancer multi-omics studies integrate complementary molecular layers such as transcriptomics, copy-number alteration, methylation, protein abundance, and mutation profiles. These data can reveal subtype structure and tumor biology that is not visible from a single assay \cite{tcga_brca_2012,hoadley_pancancer_2018}. As machine learning methods for multi-omics prediction have become increasingly complex, benchmarking has often focused on model architecture and aggregate performance \cite{wang_mogonet_2021,leng_benchmark_2022}. However, high model scores alone do not establish that a multi-omics workflow is reliable.

Three data-level issues are especially important. First, multi-omics fusion assumes that the molecular vectors being fused belong to the same biological sample. If expression from one sample is accidentally paired with copy-number, methylation, or mutation features from another sample, the marginal distribution of each modality can remain plausible while the cross-modal biological relationship is broken. Second, not all modalities have comparable sample and feature coverage. In public cancer resources, protein measurements are often sparser than transcriptomic or mutation features, and naive early fusion can turn a nominally multi-omics table into a heavily imputed table. Third, external validation cohorts can differ from the source cohort because of assay platforms, processing pipelines, enrollment criteria, and label distributions. This cohort shift can alter model behavior even when the prediction label is nominally the same.

These issues are well recognized in high-throughput biology through work on batch effects and dataset shift \cite{johnson_combat_2007,leek_batch_2010,quionero_dataset_shift_2009}, but they are not always quantified in predictive multi-omics benchmarking. The present manuscript intentionally changes the question from whether a particular architecture wins to how much data alignment, modality coverage, and cohort shift shape the reliability of multi-omics prediction.

The primary benchmark in this study is a five-task TCGA cancer-internal analysis with explicit sample-alignment and modality-coverage audits. Around that primary benchmark, we add three secondary analyses: a controlled mismatch stress test, a seven-model 10-fold model-family comparison, and a BRCA external-validation case study using source-to-external cohort-shift summaries. We also include an exploratory 10-cancer endpoint landscape screen to show that clinically interpretable within-cancer tasks are available beyond the five deeply analysed tasks. This hierarchy is intentional: the manuscript does not claim that one model architecture is universally best, nor that external validation has been completed for all cancer types. Instead, it asks how alignment, coverage, endpoint biology, and cohort shift should be made visible when cancer multi-omics prediction models are benchmarked.
